
\begin{theorem}[Positivity of bipartite partial traces]
    \label{thm:matrix_possemidef_partial_traces}
    \lean{Matrix.PosSemidef.traceLeft}
    \lean{Matrix.PosSemidef.traceRight}
    \leanok
    If $\rho$ is positive semidefinite on a bipartite space $A \otimes B$, then
    both reduced operators $\tr_A \rho$ and $\tr_B \rho$ are positive
    semidefinite.
\end{theorem}

\begin{proof}
    \leanok
    Write each partial trace as the finite sum, over the traced index, of a
    principal submatrix of $\rho$. Principal submatrices of a positive
    semidefinite matrix are positive semidefinite, and finite sums of positive
    semidefinite matrices remain positive semidefinite.
\end{proof}

\begin{definition}[Hayashi--Markov extraction of explicit $\eta$-operators]
    \label{def:mpdo_eta_of_hayashi_markov}
    \lean{MPOTensor.ExplicitEtaOperators.ofHayashiMarkov}
    \leanok
    \uses{def:mpdo_eta_structure, def:mpdo_explicit_eta_operators,
        thm:matrix_possemidef_partial_traces}
    Every Hayashi decomposition witness $h_\eta$ canonically produces an
    explicit neighboring $\eta$-family by Kronecker-multiplying the
    sector-reduced states on the neighboring bond spaces:
    \[
        \eta_{k,h}
        := \tr_C(\rho_{b_2 C}^{(k)})
        \otimes \tr_A(\rho_{A b_1}^{(h)}).
    \]
    Positive semidefiniteness of each $\eta_{k,h}$ follows from the fact
    that partial traces preserve positivity and that positive
    semidefiniteness is closed under Kronecker products. This is the
    sector-reduced extraction; the remaining connection to
    \cite[Appendix~C.2]{Cirac2016MPDO_arXiv} is the identification of this
    family with the inverse-map
    construction in the original simple-MPDO tensor coordinates.
\end{definition}

\begin{theorem}[Trace matrix of the Hayashi--Markov $\eta$-operators]
    \label{thm:mpdo_eta_of_hayashi_markov_trace_matrix}
    \lean{MPOTensor.ExplicitEtaOperators.traceMatrix_ofHayashiMarkov}
    \lean{MPOTensor.ExplicitEtaOperators.traceMatrixRe_ofHayashiMarkov}
    \lean{MPOTensor.ExplicitEtaOperators.trace_traceMatrixRe_ofHayashiMarkov}
    \leanok
    \uses{def:mpdo_eta_of_hayashi_markov,
        def:mpdo_explicit_eta_trace_matrix}
    The trace matrix induced by the extracted family is entrywise
    $T_{k,h} = 1$ in both its complex-valued form and its real-valued form.
    If there are $m$ sectors, then $T=\mathbf{1}\mathbf{1}^{\mathsf T}$ and
    $\tr(T)=m$.
\end{theorem}

\begin{proof}
    \leanok
    \uses{def:mpdo_eta_of_hayashi_markov,
        def:mpdo_explicit_eta_trace_matrix}
    Each $\rho_{b_2 C}^{(k)}$ is a density matrix, hence has unit trace; the
    partial trace $\tr_C$ preserves the trace. The analogous statement holds
    for $\tr_A \rho_{A b_1}^{(h)}$. Multiplicativity of the trace under
    Kronecker products gives $T_{k,h} = 1 \cdot 1 = 1$, and taking real parts
    gives the real-valued statement. Summing the $m$ diagonal entries gives
    $\tr(T)=m$.
\end{proof}

\begin{theorem}[Positive semidefiniteness of the Hayashi--Markov trace matrix]
    \label{thm:mpdo_eta_of_hayashi_markov_trace_matrix_psd}
    \lean{MPOTensor.ExplicitEtaOperators.traceMatrixRe_ofHayashiMarkov_posSemidef}
    \leanok
    \uses{thm:mpdo_eta_of_hayashi_markov_trace_matrix}
    The real trace matrix of the sector-reduced family is positive
    semidefinite.
\end{theorem}

\begin{proof}
    \leanok
    \uses{thm:mpdo_eta_of_hayashi_markov_trace_matrix}
    This trace matrix is the all-ones matrix
    $\mathbf{1}\mathbf{1}^{\mathsf T}$, hence is positive semidefinite.
\end{proof}

\begin{theorem}[Entrywise nonnegativity of the real trace matrix]
    \label{thm:mpdo_explicit_eta_trace_matrix_nonneg}
    \lean{MPOTensor.ExplicitEtaOperators.traceMatrixRe_nonneg}
    \leanok
    \uses{def:mpdo_explicit_eta_trace_matrix}
    Positivity of the neighboring operators gives entrywise nonnegativity of
    the real trace matrix. Matrix-level positive semidefiniteness of $T$ is a
    separate structural condition.
\end{theorem}

\begin{proof}
    \leanok
    The trace of a positive semidefinite neighboring operator is non-negative in
    the complex order. Taking real parts gives the stated entrywise
    nonnegativity.
\end{proof}

\begin{definition}[Auxiliary matrix predicates for the rank-one step]
    \label{def:matrix_trace_rankone_predicates}
    \lean{Matrix.TracePowersConstant}
    \lean{Matrix.HasRankOneFactorization}
    \leanok
    The auxiliary predicates state two matrix-theoretic ingredients used in
    the rank-one step of \cite[Appendix~C.2, Lemma~C.5]{Cirac2016MPDO_arXiv}:
    constant traces of all positive powers and existence of a rank-one
    factorization $T = a b^\top$.

    The proposed implication from primitivity and constant trace powers to a
    rank-one factorization is not a valid general theorem. A
    concrete $3 \times 3$ counterexample appears in the archive. The corrected
    criterion proved here adds positive semidefiniteness and trace
    normalization, which provide the diagonalizability needed to turn the
    trace-power condition into a rank-one factorization.
\end{definition}
% The deviation from the source is documented in
% docs/paper-gaps/cpgsv17_pf_rank_one.tex.

\begin{theorem}[Trace square of a Hermitian matrix]
    \label{thm:matrix_hermitian_trace_square_eigenvalues}
    \lean{Matrix.IsHermitian.trace_sq_eq_sum_eigenvalues_sq}
    \leanok
    Let $T$ be a finite Hermitian matrix over a real or complex scalar field,
    with Hermitian eigenvalues $\lambda_i$. Then
    \[
        \tr(T^2) = \sum_i \lambda_i^2.
    \]
\end{theorem}

\begin{proof}
    \leanok
    The spectral theorem writes $T$ as a unitary conjugate of the diagonal
    matrix of its Hermitian eigenvalues. Squaring commutes with this conjugation,
    and cyclicity of trace reduces the identity to the trace of the squared
    diagonal matrix.
\end{proof}

\begin{theorem}[Rank-one eigenvector resolution of a Hermitian matrix]
    \label{thm:matrix_hermitian_rank_one_eigenvector_resolution}
    \lean{Matrix.IsHermitian.eq_sum_eigenvalues_smul_eigenvectorProjection}
    \leanok
    Let $A$ be a finite Hermitian matrix. Choose an orthonormal eigenbasis
    $(u_k)_k$ with eigenvalues $(\lambda_k)_k$, and let
    $E_k=\ketbra{u_k}{u_k}$ be the rank-one projection onto $u_k$. Then
    \[
        A=\sum_k \lambda_k E_k.
    \]
\end{theorem}

\begin{proof}
    \leanok
    The finite-dimensional spectral theorem gives
    $A=U\diag(\lambda_k)U^\dagger$, where the columns of $U$ are
    the vectors $u_k$. Expanding the diagonal matrix as the sum of its
    rank-one coordinate projections gives the stated identity.
\end{proof}

\begin{theorem}[PSD trace-power rank-one criterion]
    \label{thm:matrix_psd_trace_power_rank_one}
    \lean{Matrix.PosSemidef.trace_powers_constant_implies_rank_one}
    \leanok
    \uses{def:matrix_trace_rankone_predicates}
    Let $T$ be a finite real square matrix. If $T$ is positive semidefinite,
    normalized by $\tr(T)=1$, and has constant traces on all positive powers,
    then $T$ has a rank-one factorization.
\end{theorem}

\begin{proof}
    \leanok
    \uses{def:matrix_trace_rankone_predicates, thm:matrix_hermitian_trace_square_eigenvalues}
    The spectral theorem diagonalizes the positive semidefinite matrix with
    non-negative real eigenvalues. The trace gives that their sum is $1$, while
    the second trace moment gives that the sum of their squares is also $1$.
    Hence exactly one eigenvalue is equal to $1$ and all others vanish, so the
    diagonal form is rank one. Unitary conjugation preserves the existence of an
    outer-product factorization.
\end{proof}

\begin{theorem}[Idempotent trace-one criterion]
    \label{thm:matrix_idempotent_trace_one_rank_one}
    \lean{Matrix.hasRankOneFactorization_of_mul_self_eq_self}
    \leanok
    Let $T$ be a finite real square matrix. If
    \[
        T^2=T,
        \qquad \tr(T)=1,
    \]
    then there are real vectors $a,b$ such that $T=a b^\top$.
\end{theorem}

\begin{proof}
    \leanok
    \uses{def:matrix_trace_rankone_predicates}
    For an idempotent endomorphism $f$ one has
    \[
        \tr(f)=\dim(\range(f)).
    \]
    Hence the range of $T$ has dimension one. Choosing a nonzero vector $a$
    spanning it expresses every column of $T$ as $b_j a$ for a scalar $b_j$,
    and therefore $T=ab^\top$.
\end{proof}

\begin{theorem}[Constant trace powers of an idempotent matrix]
    \label{thm:matrix_idempotent_trace_powers_constant}
    \lean{Matrix.tracePowersConstant_of_mul_self_eq_self}
    \leanok
    \uses{def:matrix_trace_rankone_predicates}
    Let $T$ be a finite real square matrix with $T^2 = T$. Then for every
    positive integer $N$,
    \[
        \tr(T^N) = \tr(T).
    \]
    This recovers the trace identity displayed in the proof of
    \cite[Appendix~C.2, Lemma~C.5]{Cirac2016MPDO_arXiv} from idempotence.
\end{theorem}

\begin{proof}
    \leanok
    Induction on $N$: for $N \geq 1$,
    $T^{N+1} = T^{N} T = T T = T$, so every positive power of $T$ equals $T$.
\end{proof}

\begin{theorem}[Positive entries in every row and column of a primitive matrix]
    \label{thm:matrix_primitive_row_col_pos}
    \lean{Matrix.IsPrimitive.exists_row_pos}
    \lean{Matrix.IsPrimitive.exists_col_pos}
    \leanok
    Let $T$ be a primitive non-negative matrix. Then every row of $T$ contains
    a strictly positive entry, and so does every column.
\end{theorem}

\begin{proof}
    \leanok
    If row $k$ of $T$ vanishes, then for every $m \geq 1$,
    \[
        (T^{m})_{k,k} = \sum_j T_{k,j} (T^{m-1})_{j,k} = 0,
    \]
    contradicting the strict entrywise positivity of some power of $T$.
    Likewise, if column $h$ of $T$ vanishes, then for every $m \geq 1$,
    \[
        (T^{m})_{h,h} = \sum_j (T^{m-1})_{h,j}\, T_{j,h} = 0,
    \]
    again contradicting the strict positivity of some power of $T$.
\end{proof}

\begin{theorem}[Rank of a strictly positive idempotent]
    \label{thm:matrix_positive_idempotent_rank_one}
    \lean{Matrix.rank_eq_one_of_pos_of_mul_self_eq_self}
    \leanok
    Let $I$ be a nonempty finite index set, and let $P$ be a real matrix
    indexed by $I\times I$. Suppose that $P^2=P$ and, for every $i,j\in I$,
    \[
        P_{i,j}>0.
    \]
    Then
    \[
        \rank(P)=1.
    \]
\end{theorem}

\begin{proof}
    \leanok
    Put $u=P\mathbf 1$, so every coordinate of $u$ is positive and $Pu=u$.
    For a fixed vector $x$ satisfying $Px=x$, choose an index $k$ at which
    $x_i/u_i$ is maximal and put $c=x_k/u_k$. Then $y=cu-x$ is non-negative,
    $y_k=0$, and $Py=y$. Hence
    \[
        0=y_k=(Py)_k=\sum_j P_{k,j}y_j.
    \]
    Since every $P_{k,j}$ is strictly positive, all coordinates of $y$ vanish.
    Thus every fixed vector is proportional to $u$. The range of an idempotent
    is its fixed space, so it is one-dimensional.
\end{proof}

\begin{theorem}[Rank and factorization of a stationary primitive square]
    \label{thm:matrix_primitive_pow_two_rank_one}
    \lean{Matrix.IsPrimitive.smul_of_pos}
    \lean{Matrix.IsPrimitive.rank_pow_two_eq_one_of_pow_two_eq_pow_three}
    \lean{Matrix.IsPrimitive.exists_pow_two_eq_vecMulVec_of_pow_two_eq_pow_three}
    \leanok
    Let $T$ be a primitive real $N\times N$ matrix, where
    $N\geq 1$.  Multiplication by a positive scalar preserves primitivity.  If
    \[
        T^2=T^3,
    \]
    then
    \[
        \rank(T^2)=1.
    \]
    There are vectors $a,b\in\mathbb R^N$ such that
    \[
        T^2=ab^{\mathsf T},
        \qquad
        \sum_i a_i b_i=1.
    \]
    This finite-dimensional matrix lemma supports the argument sought at
    \cite[Appendix~C.2, line~1613]{Cirac2016MPDO_arXiv}; it is not a theorem
    asserted there.
\end{theorem}

\begin{proof}
    \leanok
    \uses{thm:matrix_positive_idempotent_rank_one,
        thm:matrix_pairing_sq_eq_cube,
        thm:matrix_idempotent_trace_one_rank_one}
    Primitivity gives $m\geq 1$ such that $(T^m)_{i,j}>0$ for every $i,j$.
    Stabilization gives $T^{2m}=T^2$. Thus, for every $i,j$,
    \[
        (T^2)_{i,j}=(T^{2m})_{i,j}
        =\sum_k (T^m)_{i,k}(T^m)_{k,j}>0.
    \]
    The same stabilization identity gives $(T^2)^2=T^2$.
    The preceding theorem applied to $P=T^2$ now gives the conclusion.
    Since an idempotent has trace equal to its rank,
    $\tr(T^2)=1$.  Factoring its one-dimensional range gives
    $T^2=ab^{\mathsf T}$, and taking the trace yields
    $\sum_i a_i b_i=1$.
\end{proof}

\begin{theorem}[Idempotence of the sector trace matrix]
    \label{thm:matrix_pairing_idempotent}
    \lean{Matrix.mul_self_eq_self_of_pairing_idempotent}
    \lean{Matrix.mul_self_eq_self_of_pairing_idempotent_dual}
    \leanok
    Let $V$ be a module over a commutative ring, let $l_1,\dots,l_n \in V$,
    let $r_1,\dots,r_n$ be linear functionals on $V$, and let
    $T_{k,h}=(r_k|l_h)$ be the pairing matrix; for the tensors of
    \cite[Appendix~C.2]{Cirac2016MPDO_arXiv} this is the sector trace matrix.
    Assume the identity
    \[
        \sum_k |l_k)(r_k| = \sum_{k,h} T_{k,h}\, |l_k)(r_h|,
    \]
    whose right-hand side is the square of its left-hand side, so that the
    operator $M = \sum_k |l_k)(r_k|$ satisfies $M^2 = M$. If the vectors
    $l_1,\dots,l_n$ are linearly independent, then
    \[
        T^2 = T.
    \]
    The same conclusion holds if instead the functionals $r_1,\dots,r_n$ are
    linearly independent.
\end{theorem}

\begin{proof}
    \leanok
    Evaluating $M^2 v = M v$ at an arbitrary $v \in V$ gives
    \[
        \sum_k \bigl((r_k|Mv) - (r_k|v)\bigr)\, l_k = 0,
    \]
    so linear independence of the $l_k$ yields $(r_k|Mv)=(r_k|v)$ for all $k$
    and $v$. Taking $v = l_h$ and expanding
    $M l_h = \sum_j T_{j,h}\, l_j$ turns the left-hand side into
    $\sum_j T_{k,j} T_{j,h} = (T^2)_{k,h}$ and the right-hand side into
    $T_{k,h}$. For linearly independent functionals, apply the same argument
    in the dual space, with the $r_k$ as vectors and evaluation at $l_k$ as
    functionals.  The dual pairing matrix therefore satisfies
    \[
        (T^\top)^2=T^\top.
    \]
    Taking transposes gives
    \[
        T^2=T.
    \]
\end{proof}

Equivalently, the zero-correlation-length identity is the algebraic equality
\[
    \sum_k |l_k)(r_k|
    = \sum_{k,h} T_{k,h}\,|l_k)(r_h|,
    \qquad T_{k,h}=(r_k|l_h),
\]
which is the displayed identity of
\cite[Appendix~C.2, lines~1490--1493]{Cirac2016MPDO_arXiv}.

\begin{theorem}[Unconditional constant trace powers of the pairing matrix]
    \label{thm:matrix_pairing_sq_eq_cube}
    \lean{Matrix.pow_two_eq_pow_three_of_pairing_idempotent}
    \lean{Matrix.trace_eq_trace_sq_of_pairing_idempotent}
    \lean{Matrix.pow_eq_pow_two_of_pow_two_eq_pow_three}
    \lean{Matrix.tracePowersConstant_of_pairing_idempotent}
    \leanok
    \uses{def:matrix_trace_rankone_predicates}
    Let $V$ be a real vector space, let $l_1,\dots,l_n \in V$, let
    $r_1,\dots,r_n$ be real linear functionals on $V$, and let
    $T_{k,h}=(r_k|l_h)$ be the pairing matrix, as in
    Theorem~\ref{thm:matrix_pairing_idempotent}. Assuming only the identity
    $M^2 = M$ for the operator $M = \sum_k |l_k)(r_k|$, without assuming
    linear independence of $l_1,\dots,l_n$ or of $r_1,\dots,r_n$,
    \[
        T^2 = T^3,
    \]
    so $T^N = T^2$ for every $N \geq 2$, and for every positive integer $N$,
    \[
        \tr(T^N) = \tr(T).
    \]
\end{theorem}

\begin{proof}
    \leanok
    \uses{thm:matrix_pairing_idempotent}
    Multiplying the identity $M^2 = M$ on the left by the functional map and
    on the right by the coefficient map turns the two sides directly into
    $T^2$ and $T^3$, with no independence hypothesis; induction then gives
    $T^N = T^2$ for every $N \geq 2$. For the trace identity, restrict the
    operator $M$ to the finite-dimensional span of $l_1,\dots,l_n$: there it
    is an idempotent endomorphism of a finite free module, and cyclicity of
    the trace of a composition between that span and the coefficient space
    identifies $\tr(T)$ with the trace of the restriction and
    $\tr(T^2)$ with the trace of its square, which agree since
    the restriction is idempotent.
\end{proof}

% Scope restriction (linear independence): the source Lemma C.5 neither
% assumes nor derives linear independence of the sector tensors. Its proof
% instead uses primitivity of the trace matrix, obtained from the injectivity
% of the tensor K in the preceding construction (arXiv:1606.00608,
% lines 1457--1470), in a Perron--Frobenius trace-power argument, and
% primitivity does not entail the independence. The hypothesis below remains
% to be discharged at the matrix product density operator call site.
% Recorded in docs/paper-gaps/cpgsv17_pf_rank_one.tex.
\begin{theorem}[Rank-one factorization of the sector trace matrix]
    \label{thm:matrix_pairing_rank_one}
    \lean{Matrix.hasRankOneFactorization_of_pairing_idempotent}
    \leanok
    \uses{def:matrix_trace_rankone_predicates}
    In the setting of Theorem~\ref{thm:matrix_pairing_idempotent} with real
    scalars and linearly independent vectors $l_1,\dots,l_n$, if in addition
    $\tr(T)=1$, then there are real vectors $a,b$ such that
    $T = a b^\top$. This is the factorization $T_{k,h} = a_k b_h$ obtained in
    the proof of \cite[Appendix~C.2, Lemma~C.5]{Cirac2016MPDO_arXiv} for the
    sector trace matrix, derived here from the displayed
    zero-correlation-length identity together with the linear independence of
    the $l_k$, instead of the Perron--Frobenius trace-power argument of the
    source.
\end{theorem}

\begin{proof}
    \leanok
    \uses{thm:matrix_pairing_idempotent, thm:matrix_idempotent_trace_one_rank_one}
    Theorem~\ref{thm:matrix_pairing_idempotent} gives $T^2 = T$, and the
    idempotent trace-one criterion of
    Theorem~\ref{thm:matrix_idempotent_trace_one_rank_one} gives the
    factorization $T = a b^\top$.
\end{proof}

\begin{theorem}[PSD-corrected rank-one $T$ criterion]
    \label{thm:psd_trace_powers_rank_one_t}
    \lean{MPOTensor.sal_zcl_implies_rank_one_T_of_posSemidef}
    \leanok
    \uses{def:matrix_trace_rankone_predicates}
    Let $T$ be a finite real square matrix. Assume that $T$ is primitive,
    positive semidefinite, normalized by $\tr(T)=1$, and has constant traces on
    all positive powers. Then
    \[
        T = a b^\top
    \]
    for some vectors $a$ and $b$, with $a \cdot b = 1$.
\end{theorem}

\begin{proof}
    \leanok
    \uses{def:matrix_trace_rankone_predicates,
        thm:matrix_psd_trace_power_rank_one}
    The positive-semidefinite matrix criterion diagonalizes $T$, uses the first
    two trace moments to show that exactly one eigenvalue is nonzero, and hence
    gives the factorization. Taking traces and using
    $\tr(a b^\top)=a\cdot b$ gives the normalization $a\cdot b=1$.
\end{proof}

\begin{definition}[Sector tensors paired into the sector trace matrix]
    \label{def:mpdo_sector_pairing_data}
    \lean{MPOTensor.SectorPairingData}
    \leanok
    For a real $n \times n$ matrix $T$ and a real vector space $V$, this
    structure consists of vectors $|l_0), \dots, |l_{n-1}) \in V$ and linear
    functionals $(r_0|, \dots, (r_{n-1}|$ on $V$ satisfying the pairing identity
    \[
        T_{k,h} = (r_k|l_h)
    \]
    of \cite[Appendix~C.2]{Cirac2016MPDO_arXiv} together with the
    zero-correlation-length identity
    \[
        \sum_k |l_k)(r_k| = \sum_{k,h} T_{k,h}\, |l_k)(r_h|
    \]
    displayed in the proof of \cite[Lemma~C.5]{Cirac2016MPDO_arXiv}. The
    vectors are the closed sector tensors obtained there from the sector
    splitting of an injective simple tensor.
\end{definition}

\begin{theorem}[Idempotence of the sector trace matrix from the zero-correlation-length identity]
    \label{thm:mpdo_sector_pairing_idempotent}
    \lean{MPOTensor.SectorPairingData.pairing_operator_idempotent}
    \lean{MPOTensor.SectorPairingData.mul_self_eq_self}
    \leanok
    \uses{def:mpdo_sector_pairing_data}
    For sector tensors paired into $T$ as in
    Definition~\ref{def:mpdo_sector_pairing_data}, the operator
    $M = \sum_k |l_k)(r_k|$ satisfies $M^2 = M$. If in addition the vectors
    $|l_0), \dots, |l_{n-1})$ are linearly independent, then $T^2 = T$.
\end{theorem}

\begin{proof}
    \leanok
    \uses{thm:matrix_pairing_idempotent}
    Substituting $T_{k,h} = (r_k|l_h)$ into the right-hand side of the
    zero-correlation-length identity gives, for every $v \in V$,
    \[
        \sum_{k,h} T_{k,h}(r_h|v)\,|l_k)
        = \sum_{k,h} (r_k|l_h)(r_h|v)\,|l_k)
        = \sum_k (r_k|Mv)\,|l_k)
        = M^2 v,
    \]
    so $M^2 = M$. Linear independence of the $|l_k)$ then
    gives $T^2 = T$ by Theorem~\ref{thm:matrix_pairing_idempotent}.
\end{proof}

\begin{theorem}[Unconditional constant trace powers of the sector trace matrix]
    \label{thm:mpdo_sector_pairing_trace_powers}
    \lean{MPOTensor.SectorPairingData.pairing_sq_eq_pairing_cube}
    \lean{MPOTensor.SectorPairingData.tracePowersConstant}
    \leanok
    \uses{def:mpdo_sector_pairing_data}
    For sector tensors paired into $T$ as in
    Definition~\ref{def:mpdo_sector_pairing_data}, without any independence
    hypothesis on the closed sector tensors or the functionals,
    \[
        T^2 = T^3,
    \]
    and for every positive integer $N$,
    \[
        \tr(T^N) = \tr(T).
    \]
\end{theorem}

\begin{proof}
    \leanok
    \uses{thm:mpdo_sector_pairing_idempotent, thm:matrix_pairing_sq_eq_cube}
    The operator identity $M^2 = M$ established in the proof of
    Theorem~\ref{thm:mpdo_sector_pairing_idempotent} gives $T^2 = T^3$ and
    the constant trace powers by Theorem~\ref{thm:matrix_pairing_sq_eq_cube},
    with no independence hypothesis.
\end{proof}

% Scope restriction (independent subspaces): the independent subspace family in the
% second statement is not displayed in arXiv:1606.00608; the closed sector
% tensors of that construction live in a common bond space after the physical
% sector legs are contracted. Whether the inverse-map construction supplies
% such supports, or the linear independence by other means, remains open.
% Recorded in docs/paper-gaps/cpgsv17_pf_rank_one.tex.
\begin{theorem}[Nonvanishing and independence of the sector tensors]
    \label{thm:mpdo_sector_tensors_independent}
    \lean{MPOTensor.SectorPairingData.l_ne_zero}
    \lean{MPOTensor.SectorPairingData.linearIndependent_l}
    \leanok
    \uses{def:mpdo_sector_pairing_data, thm:matrix_primitive_row_col_pos}
    Suppose the sector trace matrix $T$ of
    Definition~\ref{def:mpdo_sector_pairing_data} is primitive. Then every
    closed sector tensor satisfies $|l_k) \neq 0$. If, in addition, each
    $|l_k)$ lies in its own member of an independent family of subspaces of
    $V$, then
    the vectors $|l_0), \dots, |l_{n-1})$ are linearly independent.
\end{theorem}

\begin{proof}
    \leanok
    \uses{thm:matrix_primitive_row_col_pos}
    By Theorem~\ref{thm:matrix_primitive_row_col_pos}, some entry of the
    $k$-th column of $T$ is positive, and that entry is the pairing
    $(r_j|l_k)$, so $|l_k) \neq 0$. Nonzero vectors chosen from an independent
    family of subspaces, one from each member, are linearly independent.
\end{proof}

\begin{theorem}[Rank one from sector-pairing idempotence under linear independence]
    \label{thm:pairing_idempotent_rank_one_t}
    \lean{MPOTensor.sal_zcl_implies_rank_one_T_of_pairing_idempotent}
    \lean{MPOTensor.sal_zcl_implies_rank_one_T_of_sector_supports}
    \leanok
    \uses{def:mpdo_sector_pairing_data, def:matrix_trace_rankone_predicates}
    Let $T$ be a sector trace matrix with $\tr(T)=1$, paired
    from sector tensors satisfying the zero-correlation-length identity of
    Definition~\ref{def:mpdo_sector_pairing_data}. If the closed sector
    tensors $|l_0), \dots, |l_{n-1})$ are linearly independent, then
    \[
        T_{k,h} = a_k b_h,
        \qquad
        \sum_k a_k b_k = 1
    \]
    for some real numbers $a_k, b_h$. The same conclusion holds if, instead
    of the independence, $T$ is primitive and the closed sector tensors lie
    in distinct members of an independent family of subspaces. In contrast
    with the positive-semidefinite criterion of
    Theorem~\ref{thm:psd_trace_powers_rank_one_t}, the constant trace powers
    are derived rather than assumed.
\end{theorem}

\begin{proof}
    \leanok
    \uses{thm:mpdo_sector_pairing_idempotent,
        thm:mpdo_sector_pairing_trace_powers,
        thm:mpdo_sector_tensors_independent,
        thm:matrix_idempotent_trace_one_rank_one}
    Theorem~\ref{thm:mpdo_sector_pairing_idempotent} gives $T^2 = T$, and
    Theorem~\ref{thm:mpdo_sector_pairing_trace_powers} gives the constant
    trace powers unconditionally. The idempotent trace-one criterion of
    Theorem~\ref{thm:matrix_idempotent_trace_one_rank_one} yields the
    factorization $T = a b^\top$, and taking traces gives
    $\sum_k a_k b_k = \tr(T) = 1$. When the independence is not
    assumed directly, Theorem~\ref{thm:mpdo_sector_tensors_independent}
    derives it from primitivity once each sector tensor lies in its own member
    of an independent family of subspaces.
\end{proof}
